% !PRINTED 688
\weberpage{688}
Denken wir uns das Gebiet $S$ gegeben, so wird das entsprechende
 Gebiet $S'$ von $t$ abhängig sein und sich mit $t$ vergrössern.
Die Punkte mit \gesperrt{ganzzahligen} Coordinaten $x_i$ nennen
wir, wie früher, die \gesperrt{Gitterpunkte}. Die Anzahl der in einem
endlichen Gebiete $S'$ liegenden Gitterpunkte ist immer endlich,
wächst aber mit $t$ und soll mit $Z_t$ bezeichnet werden. Dann ist
nach den Sätzen \S.\,155, (14) das Volumen des Gebietes $S$, d.\,h. das
über alle Punkte von $S$ erstreckte $n$-fache Integral
\[
V = \int\!\!\int \ldots \int dx_1\, dx_2 \ldots dx_n
\]
gleich dem Grenzwerthe des Verhältnisses $Z_t : t$ für unendlich
wachsendes $t$:
\weq{5}{
V = \mathop{\mathrm{Lim}}\limits_{t=\infty} \frac{Z_t}{t}.
}
Hierin liegt zugleich die genauere Präcisirung, die wir über
das Gebiet $S$ machen müssen, wenn unsere Betrachtung anwendbar
 sein soll. Es muss $S$ so beschaffen sein, dass das
$n$-fache Integral $V$ eine bestimmte Bedeutung hat. Jeder Gitterpunkt
 ist nun durch (1) das Bild einer ganzen durch $\mathfrak{a}$ theilbaren
 Zahl des Körpers $\Omega$, und wir können also auch sagen,
dass $Z_t$ die Anzahl der durch $\mathfrak{a}$ theilbaren ganzen Zahlen in $\Omega$
ist, deren Bilder in dem Gebiete $S'$ liegen.

Um das Gebiet $S'$ in geeigneter Weise abzugrenzen, wählen
wir ein System fundamentaler Einheiten des Körpers $\Omega$
\[
\varepsilon_1, \varepsilon_2, \ldots, \varepsilon_{\nu-1}
\]
mit dem System der conjugirten Logarithmen
\[
l_{i,1}, l_{i,2}, \ldots, l_{i,\nu}
\]
und dem Regulator
\weq{6}{
\pm L(\varepsilon_1, \varepsilon_2, \ldots, \varepsilon_{\nu-1}) = \pm \Sigma \pm l_{1,1} l_{2,2} \ldots l_{\nu-1,\nu-1},
}
der mit $L$ bezeichnet werden mag.

Wir definiren jetzt ein System von Functionen $\xi_1, \xi_2, \ldots, \xi_\nu$
der Variablen $x_i$ durch folgende Bedingungen:

Es sei, wenn die $y_i$ die Bedeutung (2) haben,
\weq{7}{
z_i = \delta_i \log |y_i|,
}
so dass, wenn für die $x_i$ ein System ganzer rationaler Zahlen
gesetzt wird, $z_i$ nach \S.\,185 in das System der conjugirten
Logarithmen der Zahl $\alpha$ übergeht.

Es ist dann, wenn $y_k$ zu einem reellen Körper gehört,
\weq{8}{
z_k = \tfrac{1}{2} \log y_k^2,
}
%CONT>>
